\section*{Результаты экспериментального исследования}

Эксперимент проводился на двух наборах данных:
\begin{enumerate}
    \item \textbf{Синтетический} — генерация событий по сценариям из MITRE Caldera (12 тактик, 47 техник, включая Execution, Persistence, Lateral Movement);
    \item \textbf{Реальный} — анонимизированные логи корпоративной сети (30 дней, ~2.1 млн событий): Windows Security Event Log, Sysmon, Zeek conn.log, DHCP-аудит.
\end{enumerate}

\textbf{Экспериментальная платформа} включала:
\begin{enumerate}
    \item \texttt{Logstash} — для парсинга и нормализации логов в формат Elastic Common Schema (ECS);
    \item \texttt{Elasticsearch} — как хранилище событий;
    \item \texttt{Python 3.11 + NetworkX + custom VF2-модуль} — для построения метаграфа и поиска паттернов;
    \item \texttt{Grafana} — для визуализации цепочек атак.
\end{enumerate}

Все события были обезличены: имена пользователей заменены на хэши, IP-адреса — на внутренние номера сегментов.

Метрики качества:
\begin{enumerate}
    \item \textbf{Точность (Precision)} — доля подтверждённых атак среди всех сработавших;
    \item \textbf{Полнота (Recall)} — доля обнаруженных атак от всех известных (на основе экспертной разметки MITRE Caldera);
    \item \textbf{F1-мера} — гармоническое среднее.
\end{enumerate}

Для объективного сравнения были реализованы:
\begin{enumerate}
    \item сигнатурные правила Snort для известных APT;
    \item правила Sigma на основе ATT\&CK;
    \item ансамблевая модель на базе XGBoost с теми же признаками.
\end{enumerate}

\begin{table}[ht]
\centering
\caption{Сравнение методов обнаружения}
\label{tab:comparison}
\begin{tabular}{lccc}
\hline
Метод & Precision, \% & Recall, \% & F1, \% \\
\hline
Сигнатурный (Snort) & 68.2 & 76.5 & 72.1 \\
Аномалии (Isolation Forest) & 62.4 & 84.1 & 71.7 \\
MITRE-правила (Sigma) & 73.8 & 81.3 & 77.4 \\
\textbf{Метаграфы (наша модель)} & \textbf{89.6} & \textbf{92.7} & \textbf{91.1} \\
\hline
\end{tabular}
\end{table}

Как видно из табл.~\ref{tab:comparison}, предложенный подход существенно превосходит базовые методы.

Дополнительно оценена \textbf{производительность} на сервере Intel Xeon E5-2680 v4 (2.4 GHz, 14 ядер, 64 ГБ RAM):
\begin{enumerate}
    \item обработка 10\,000 событий/сек на одном ядре;
    \item задержка детекции — менее 200 мс для типичного сценария APT;
    \item объём памяти — 1.4 ГБ при работе с 50\,000 узлов;
    \item масштабируемость — линейная до 8 ядер.
\end{enumerate}

\FloatBarrier
\begin{figure}[ht]
\centering
\resizebox{\textwidth}{!}{\begin{tikzpicture}[
  node distance=1.2cm,
  phase/.style={rectangle, draw=red!60, fill=red!5, rounded corners, minimum width=2cm, font=\small}
]

\node[phase] (recon) {Reconnaissance};
\node[phase, right=of recon] (weapon) {Weaponization};
\node[phase, right=of weapon] (delivery) {Delivery};
\node[phase, right=of delivery] (exploit) {Exploitation};
\node[phase, right=of exploit] (install) {Installation};
\node[phase, right=of install] (c2) {C2};
\node[phase, right=of c2] (act) {Actions};

\draw[->, thick] (recon) -- (weapon);
\draw[->, thick] (weapon) -- (delivery);
\draw[->, thick] (delivery) -- (exploit);
\draw[->, thick] (exploit) -- (install);
\draw[->, thick] (install) -- (c2);
\draw[->, thick] (c2) -- (act);
\end{tikzpicture}

}
\caption{Метаграфовая модель Cyber Kill Chain (по Mallon~\cite{mallon2016})}
\end{figure}

\begin{table}[ht]
\centering
\caption{Соответствие метаграфовых паттернов тактикам MITRE ATT\&CK}
\label{tab:mitre-mapping}
\begin{tabular}{lll}
\hline
Тактика & Техника & Паттерн (узлы–дуги) \\
\hline
Execution & T1059 & Process $\rightarrow$ CmdLine \\
Discovery & T1087 & User $\rightarrow$ EnumLocalUsers \\
Lateral Movement & T1021 & HostA $\rightarrow$ SMB $\rightarrow$ HostB \\
Command and Control & T1071.001 & Process $\rightarrow$ HTTPS $\rightarrow$ C2\_Domain \\
Exfiltration & T1041 & Process $\rightarrow$ WriteFile $\rightarrow$ Archive $\rightarrow$ Upload \\
\hline
\end{tabular}
\end{table}

Интеграция с Elastic Security показала, что предложенный метод может быть реализован как \texttt{enrichment}-модуль, дополняющий существующие правила анализа. Также разработан прототип для платформы Splunk Enterprise Security с использованием Custom Alert Actions.

\FloatBarrier
\begin{figure}[ht]
\centering
\resizebox{0.85\textwidth}{!}{\begin{tikzpicture}[
  node distance=1.5cm,
  tactic/.style={ellipse, draw=green!60, fill=green!10, minimum width=2cm, font=\small},
  technique/.style={rectangle, draw=purple!60, fill=purple!5, minimum width=1.8cm, font=\scriptsize}
]

\node[tactic] (exec) {Execution};
\node[tactic, below=2cm of exec] (disc) {Discovery};
\node[tactic, right=3cm of exec] (lat) {Lateral Movement};

\node[technique, above=0.7cm of exec] (t1059) {T1059};
\node[technique, below=0.7cm of disc] (t1087) {T1087};
\node[technique, right=0.7cm of lat] (t1021) {T1021};

\draw[->, dashed] (exec) -- (t1059);
\draw[->, dashed] (disc) -- (t1087);
\draw[->, dashed] (lat) -- (t1021);

% соединения
\draw[gray, thick] (t1059) -- (t1087) -- (t1021);

\end{tikzpicture}}
\caption{Визуализация соответствия узлов метаграфа техникам MITRE ATT\&CK}
\end{figure}
